\section{数一几何入口：投影曲线、柱面与积分区域}\label{r2:projection}
空间曲线 $\Gamma$ 在 $xy$ 面上的投影是所有使“存在 $z$ 满足原方程与约束”的 $(x,y,0)$。若消元得到 $g(x,y)=0$，则把 $z$ 留作任意值的 $g(x,y)=0$ 描述的是投影柱面；与 $z=0$ 联立才是平面投影曲线。立体区域的投影通常是不等式区域，而非只是一条边界曲线。

\begin{method}{四步定投影与积分限}
先写出全部曲面、方位与不等式约束；选择投影方向并消去对应变量；检查是否存在实数解及所选分支，得到投影区域；最后在每个投影点上确定上、下边界或各曲面支。消元只给必要关系时，要代回验证，不能把增根或错支当成真实投影。
\end{method}

\begin{example}{球面与平面交线的三种对象}{数一训练}
交线满足 $x^2+y^2+z^2=1$、$z=x$。消元得到 $2x^2+y^2=1$。
\begin{itemize}
\item 投影曲线：$2x^2+y^2=1,z=0$，是 $xy$ 面上的椭圆。
\item 投影柱面：$2x^2+y^2=1$，母线平行 $z$ 轴。
\item 原交线：必须保留 $z=x$，不能用投影曲线替代空间位置。
\end{itemize}
原曲线可以参数化为 $x=\cos t/\sqrt2,y=\sin t,z=\cos t/\sqrt2$，说明每个椭圆投影点都可提升回原曲线；若再限第一卦限，还需 $x,y,z\ge0$，不能保留整个椭圆。
\end{example}

\begin{example}{从存在性得到立体投影区域}{数一训练}
立体 $\Omega$ 由 $x^2+y^2\le z\le2$ 围成。存在合法 $z$ 的充要条件为 $x^2+y^2\le2$，故
\[
 D_{xy}=\{x^2+y^2\le2\},\qquad
 \iiint_\Omega f\dd V=\iint_{D_{xy}}\int_{x^2+y^2}^2 f(x,y,z)\dd z\dd A.
\]
对 $f=1$，极坐标算出体积 $2\pi\int_0^{\sqrt2}(2-r^2)r\dd r=2\pi$。若投影到 $xz$ 面，则 $0\le z\le2,|x|\le\sqrt z$，且 $-\sqrt{z-x^2}\le y\le\sqrt{z-x^2}$；两种积分次序描述同一立体。
\end{example}

\paragraph{曲面积分别忘记多层与方向}
图形面 $z=g(x,y)$ 的第一类面积元为 $\sqrt{1+g_x^2+g_y^2}\dd A$；上侧向量面积元为 $(-g_x,-g_y,1)\dd A$。球面投影到圆盘时有上、下两层，第一类积分须逐层累加，第二类须分别按方向写符号。投影边界上的 $F_z=0$ 也提示不能在整张曲面上不分支地使用 $z=z(x,y)$。
自测：$x^2+y^2+z^2=4$ 与 $z=1$ 的交线投影。答：投影曲线为 $x^2+y^2=3,z=0$；投影柱面为 $x^2+y^2=3$；原空间交线仍位于 $z=1$。
